\section{Conclusion}

Based on questionnaire data from 32 female college students and 3 semi-structured interviews, this study examined the relationships among gym spatial environment, body image, and resistance training avoidance in a Chinese/Asian campus context.
The results showed that participants were not lacking in exercise interest: the majority had used the campus gym and maintained a regular exercise frequency.
However, utilization of the free-weight area was notably low, with 68.75\% of the sample actively avoiding the free-weight area at least ``sometimes.''
Participant-level data further revealed a relatively strong positive association between spatial pressure/training anxiety and free-weight area avoidance frequency, whereas the relationship between social media appearance internalization and avoidance was weaker.
The most prominent barriers were not concern about building visible muscle, but rather the male-dominated environment, unfamiliarity with equipment, not knowing where to begin, fear of incorrect form, and the discomfort caused by open, visible spaces.

Social media appearance content may influence exercise goal selection for some individuals ($\rho = .31$, $p = .084$), but the present sample did not exhibit a strong anxiety that ``strength training will undermine a slender feminine image.''
Interviews also revealed that respondents frequently held reserved or even dismissive attitudes toward relevant stereotypes.
Therefore, the central conclusion of this paper is: in the Chinese/Asian campus context of the present sample, resistance training avoidance among female college students in campus gyms is more primarily a spatial and beginner barrier issue rather than solely an appearance internalization issue.

Regarding RQ3, this study originally hypothesized that participants would negotiate between health motivation and digital appearance norms, but the actual data more consistently pointed to a different form of negotiation: participants primarily sought strategic ways to participate within a space where they felt overly visible and lacked entry pathways, rather than making trade-offs between health and appearance. This shift is itself consistent with the preceding findings---in the present sample, spatial visibility and beginner barrier constituted more salient obstacles than appearance pressure, and the core object of ``negotiation'' was spatial conditions rather than appearance norms.

In practice, universities should prioritize beginner-friendly women's strength training workshops, equipment instruction, companion support, and semi-private/beginner-friendly spatial modifications.
These measures can transform ``having a gym'' into genuine accessibility characterized by ``daring to enter, knowing how to use equipment, and being able to sustain practice.''

Theoretically, this study demonstrates the value of analyzing spatial production, body objectification, and behavioral choices within a single framework. When ``the space does not belong to me'' and ``my body must not deviate from ideal femininity'' co-occur, spatial visibility and beginner barrier constitute more prominent obstacles than appearance internalization in the present sample. This suggests that future research should not understand women's resistance training avoidance solely from the single perspective of body image or exercise participation, but should attend to how gym spaces are experienced as social fields with entry barriers.

Given the limited sample size of this study and the fact that all individual-level statistics are exploratory bivariate analyses, the conclusions should not be extrapolated as causal laws; however, the findings clearly suggest that gender equity in campus gyms depends not only on equipment availability, but also on how spaces are experienced, who is able to make mistakes and learn within them, and whether female beginners can be seen as legitimate users.
